\section{Related Work}
\label{sec:related}

\paragraph{LLM Agents for Interactive Environments.}
The ReAct framework~\citep{yao2023react} established the paradigm of interleaving chain-of-thought reasoning with environment actions, demonstrating strong performance on knowledge-intensive and decision-making tasks. In embodied settings, ALFWorld~\citep{shridhar2021alfworld} provides a text-based household environment where agents must manipulate objects to complete tasks. WebShop~\citep{yao2022webshop} tests agents on product search and purchase in a simulated e-commerce website. ScienceWorld~\citep{wang2022scienceworld} challenges agents with multi-step science experiments requiring domain knowledge. These benchmarks share a common pattern: agents must execute sequences of actions where early mistakes cascade into failure, making error recovery a critical capability.

\paragraph{Experiential Memory for LLM Agents.}
Following the taxonomy of~\citet{liu2025memory}, experiential memory for agents can be categorized into case-based (storing raw trajectories), strategy-based (storing insights and rules), and skill-based (storing reusable action sequences). Reflexion~\citep{shinn2023reflexion} generates free-text self-reflections after failed episodes, injecting them as additional context in subsequent trials---a form of strategy-based memory with episode-level injection. ExpeL~\citep{zhao2024expel} extracts cross-task experiential insights from trajectory pools, enabling knowledge transfer across different tasks. R2D2~\citep{zheng2024r2d2} combines reasoning and retrieval for decision-making using distilled knowledge. REMEMBERER~\citep{majumder2023rememberer} augments agents with a long-term experience memory retrieved via natural language queries. All these approaches share a common design: memory is \emph{injected at episode boundaries}---at the start of a new trial or when switching tasks. None supports mid-episode retrieval triggered by real-time failure detection.

\paragraph{Memory Retrieval Timing.}
\citet{liu2025memory} distinguish between passive retrieval (triggered at fixed intervals or every step) and active retrieval (triggered by specific conditions). Most existing agent memory systems use passive timing: Reflexion injects at episode start, ExpeL retrieves before each task. Active retrieval based on internal signals---such as detecting that an action has failed---remains underexplored. Our work fills this gap by proposing failure-triggered retrieval: the agent monitors its execution for failure patterns and retrieves relevant memories only when a failure is detected, avoiding unnecessary retrieval overhead while enabling timely course correction.
